DreamDiffusion \cite{dreamdiffusion} a été le premier modèle de diffusion à
apprendre à générer des images à partir d'électroencéphalogrammes (EEG). Le
stage était axé sur la génération d'images conditionnée par les EEG, en
reproduisant les résultats de DreamDiffusion et en explorant d'autres modèles
accomplissant la même tâche mais avec des techniques nouvelles. En effet,
DreamDiffusion produisait parfois des résultats incohérents ou semblait
conditionné par la classe de l'image de référence plutôt que par ses éléments
réels — une fidélité sémantique manquait, par exemple lorsqu'il générait une
photo de chaise, mais d'un type très différent et dans un contexte tout autre.
Une tentative a été faite pour proposer une nouvelle métrique destinée aux
modèles de génération d'images à partir d'EEG, permettant de mieux mesurer la
justesse sémantique de l'image générée. L'idée consiste à comparer l'image
générée à l'image de référence en mesurant la cosine similarity entre les
embeddings (issus de sBERT \cite{sbert}) des captioning (générées par BLIP2
\cite{blip2} ou MiniCPM \cite{minicpm}) des deux images.  Ce travail n'a pas été
achevé à la fin du stage mais a été poursuivi par de nouveaux stagiaires, et
nous avons soumis un article à une conférence ne pouvant être nommée tant que le
papier n'a pas été revu. 

DreamDiffusion \cite{dreamdiffusion} was the first diffusion model that has
learnt to generate images from Electroencephalograms (EEG). The internship
focused on image generation conditioned by EEGs, replicating DreamDiffusion's
results, going through other models that also perform the same task but with
novel techniques.  Indeed, DreamDiffusion gave results that could sometimes be
incoherent, or seemed to be conditioned by the class the ground truth image was
from instead of the real elements of the ground truth, a semantic faithfulness
of the image was lacking when it generated a photo of a chair but a very
different chair in a very different context. An attempt was made at proposing a
new metric for EEG to image generation models that better measured the semantic
correctness of the generated image. The idea is to compare the generated image
to the ground truth by comparing (cosine similarity) the embedding (from sBERT
\cite{sbert}) of the caption (generated from BLIP2\cite{blip2} or
MiniCPM\cite{minicpm}) of the images. This work was not finished by the end of
the internship but continued later by new interns, and we submitted a paper to
a conference which cannot be cited as the paper is not reviewed yet.